\section{Probabilidades Condicionais}

A exposição desta parte segue de perto \cite{GrimmettWelsh2014}, que
recomendamos como leitura complementar destas notas. Introduzimos aqui as
noções de probabilidade condicional e de independência de eventos, e provamos
o teorema da partição, do qual decorre o teorema de Bayes.

Seja \(\mathcal{E}\) um experimento com espaço de probabilidade \((\Omega,
\mathcal{F}, \mathbb{P})\). Podemos, às vezes, possuir alguma informação
incompleta sobre o resultado efetivo de \(\mathcal{E}\), sem conhecer esse
resultado por inteiro. Por exemplo, se lançamos um dado honesto e um amigo nos
diz que um número par está visível, então essa informação afeta nossos cálculos
de probabilidades. Em geral, se \(A\) e \(B\) são eventos (isto é, \(A, B \in
\mathcal{F}\)) e nos é dado que \(B\) ocorre, então, à luz dessa informação, a
nova probabilidade de \(A\) pode não ser mais \(\mathbb{P}(A)\). Claramente,
nessa nova circunstância, \(A\) ocorre se, e somente se, \(A \cap B\) ocorre, o
que sugere que a nova probabilidade de \(A\) deve ser proporcional a
\(\mathbb{P}(A \cap B)\). Tornamos essa discussão mais formal em uma
definição.\footnote{Enfatizamos que isto é uma definição, e não um teorema.}

\begin{definition}\label{gw:def:prob-cond}
  Se \(A, B \in \mathcal{F}\) e \(\mathbb{P}(B) > 0\), a probabilidade
  (condicional) de \(A\) dado \(B\) é denotada por \(\mathbb{P}(A \mid B)\) e
  definida por
  \[
    \mathbb{P}(A \mid B) = \frac{\mathbb{P}(A \cap B)}{\mathbb{P}(B)}.
    \label{eq:gw:prob-cond}
  \]
\end{definition}

Observe que a constante de proporcionalidade em \eqref{eq:gw:prob-cond} foi
escolhida de modo que a probabilidade \(\mathbb{P}(B \mid B)\), de que \(B\)
ocorra dado que \(B\) ocorre, satisfaça \(\mathbb{P}(B \mid B) = 1\). Devemos,
em seguida, verificar que essa definição dá origem a um espaço de
probabilidade.

\begin{theorem}\label{gw:teo:cond-espaco-prob}
  Se \(B \in \mathcal{F}\) e \(\mathbb{P}(B) > 0\), então \((\Omega,
  \mathcal{F}, \mathbb{Q})\) é um espaço de probabilidade, onde \(\mathbb{Q}:
  \mathcal{F} \to \mathbb{R}\) é definida por \(\mathbb{Q}(A) = \mathbb{P}(A
  \mid B)\).
\end{theorem}

\begin{proof}
  Precisamos apenas verificar que \(\mathbb{Q}\) é uma medida de probabilidade
  em \((\Omega, \mathcal{F})\). Certamente \(\mathbb{Q}(A) \geq 0\) para \(A
  \in \mathcal{F}\) e
  \[
    \mathbb{Q}(\Omega) = \mathbb{P}(\Omega \mid B) =
    \frac{\mathbb{P}(\Omega \cap B)}{\mathbb{P}(B)} = 1,
  \]
  e resta verificar que \(\mathbb{Q}\) satisfaz a aditividade enumerável.
  Suponha que \(A_1, A_2, \ldots\) sejam eventos disjuntos em
  \(\mathcal{F}\). Então
  \[
    \begin{aligned}
      \mathbb{Q}\left(\bigcup_i A_i\right) &= \frac{1}{\mathbb{P}(B)}
      \mathbb{P}\left(\left(\bigcup_i A_i\right) \cap B\right) \\
      &= \frac{1}{\mathbb{P}(B)} \mathbb{P}\left(\bigcup_i (A_i \cap
        B)\right) \\
      &= \frac{1}{\mathbb{P}(B)} \sum_i \mathbb{P}(A_i \cap B) \quad
        \text{pois } \mathbb{P} \text{ satisfaz a aditividade enumerável} \\
      &= \sum_i \mathbb{Q}(A_i).
    \end{aligned}
  \]
\end{proof}

\begin{exercise}\label{gw:exercicio:regra-cadeia}
  Se \((\Omega, \mathcal{F}, \mathbb{P})\) é um espaço de probabilidade e
  \(A, B, C\) são eventos, mostre que
  \[
    \mathbb{P}(A \cap B \cap C) = \mathbb{P}(A \mid B \cap
    C)\mathbb{P}(B \mid C)\mathbb{P}(C)
  \]
  desde que \(\mathbb{P}(B \cap C) > 0\).
\end{exercise}

\begin{exercise}\label{gw:exercicio:bayes-formula}
  Mostre que
  \[
    \mathbb{P}(B \mid A) = \mathbb{P}(A \mid B)
    \frac{\mathbb{P}(B)}{\mathbb{P}(A)}
  \]
  se \(\mathbb{P}(A) > 0\) e \(\mathbb{P}(B) > 0\).
\end{exercise}

\begin{exercise}\label{gw:exercicio:sete-moedas}
  Considere o experimento de lançar uma moeda honesta 7 vezes. Determine a
  probabilidade de obter um número primo de caras, dado que cara ocorre em
  pelo menos 6 dos lançamentos.
\end{exercise}

\section{Eventos Independentes}

Dizemos que dois eventos \(A\) e \(B\) são ``independentes'' se a ocorrência
de um deles não afeta a probabilidade de o outro ocorrer. Mais formalmente,
isto exige que, se \(\mathbb{P}(A), \mathbb{P}(B) > 0\), então
\[
  \mathbb{P}(A \mid B) = \mathbb{P}(A) \quad \text{e} \quad \mathbb{P}(B
  \mid A) = \mathbb{P}(B). \label{eq:gw:indep-cond}
\]

Escrevendo \(\mathbb{P}(A \mid B) = \mathbb{P}(A \cap B)/\mathbb{P}(B)\), vemos
que a seguinte definição é apropriada.

\begin{definition}\label{gw:def:independencia}
  \emph{Eventos \(A\) e \(B\) de um espaço de probabilidade \((\Omega,
  \mathcal{F}, \mathbb{P})\) são chamados} \textbf{independentes} \emph{se}
  \[
    \mathbb{P}(A \cap B) = \mathbb{P}(A)\mathbb{P}(B), \label{eq:gw:independencia}
  \]
  \emph{e} \textbf{dependentes} \emph{caso contrário.}
\end{definition}

Esta definição é ligeiramente mais geral do que \eqref{eq:gw:indep-cond}, pois
permite que os eventos \(A\) e \(B\) tenham probabilidade zero. Ela é
facilmente generalizada, como segue, para mais de dois eventos. Uma família
\(\mathcal{A} = (A_i : i \in I)\) de eventos é chamada \emph{independente} se,
para todos os subconjuntos \emph{finitos} \(J\) de \(I\),
\[
  \mathbb{P}\left(\bigcap_{i \in J} A_i\right) = \prod_{i \in J}
  \mathbb{P}(A_i). \label{eq:gw:familia-indep}
\]
A família \(\mathcal{A}\) é chamada \emph{independente dois a dois} se
\eqref{eq:gw:familia-indep} vale sempre que \(|J| = 2\).

Assim, três eventos \(A, B, C\) são independentes se, e somente se, as
seguintes igualdades valem:
\[
  \mathbb{P}(A \cap B \cap C) = \mathbb{P}(A)\mathbb{P}(B)\mathbb{P}(C),
  \quad \mathbb{P}(A \cap B) = \mathbb{P}(A)\mathbb{P}(B),
\]
\[
  \mathbb{P}(A \cap C) = \mathbb{P}(A)\mathbb{P}(C), \quad \mathbb{P}(B \cap
  C) = \mathbb{P}(B)\mathbb{P}(C).
\]

Existem famílias de eventos que são independentes dois a dois, mas não são
independentes.

\begin{example}\label{gw:exemplo:dado-quatro-faces}
  Suponha que lançamos um dado honesto de quatro faces (você pode pensar nisto
  como um dado quadrado lançado em um universo bidimensional). Podemos tomar
  \(\Omega = \{1, 2, 3, 4\}\), onde cada \(\omega \in \Omega\) tem a mesma
  chance de ocorrer. Os eventos \(A = \{1, 2\}\), \(B = \{1, 3\}\), \(C =
  \{1, 4\}\) são independentes dois a dois, mas não são independentes.
\end{example}

\begin{exercise}\label{gw:exercicio:indep-simetrica}
  Sejam \(A\) e \(B\) eventos satisfazendo \(\mathbb{P}(A), \mathbb{P}(B) >
  0\), e tais que \(\mathbb{P}(A \mid B) = \mathbb{P}(A)\). Mostre que
  \(\mathbb{P}(B \mid A) = \mathbb{P}(B)\).
\end{exercise}

\begin{exercise}\label{gw:exercicio:disjuntos-indep}
  Se \(A\) e \(B\) são eventos disjuntos e independentes, o que se pode dizer
  sobre as probabilidades de \(A\) e \(B\)?
\end{exercise}

\begin{exercise}\label{gw:exercicio:indep-complementar}
  Mostre que os eventos \(A\) e \(B\) são independentes se, e somente se, \(A\)
  e \(\Omega \setminus B\) são independentes.
\end{exercise}

\begin{exercise}\label{gw:exercicio:indep-complementares}
  Mostre que os eventos \(A_1, A_2, \ldots, A_m\) são independentes se, e
  somente se, \(\Omega \setminus A_1, \Omega \setminus A_2, \ldots, \Omega
  \setminus A_m\) são independentes.
\end{exercise}

\begin{exercise}\label{gw:exercicio:p-eventos}
  Se \(A_1, A_2, \ldots, A_m\) são independentes e \(\mathbb{P}(A_i) = p\)
  para \(i = 1, 2, \ldots, m\), determine a probabilidade de que
  \begin{enumerate}[label=(\alph*)]
    \item nenhum dos \(A_i\) ocorra,
    \item um número par de \(A_i\) ocorra.
  \end{enumerate}
\end{exercise}

\begin{exercise}\label{gw:exercicio:dado-primo}
  Em sua mesa, há um dado muito especial com um número primo \(p\) de faces, e
  você lança este dado uma vez. Mostre que dois eventos \(A\) e \(B\) quaisquer
  não podem ser independentes, a menos que \(A\) ou \(B\) seja o espaço
  amostral inteiro ou o conjunto vazio.
\end{exercise}

\section{O Teorema da Partição}

Seja \((\Omega, \mathcal{F}, \mathbb{P})\) um espaço de probabilidade. Uma
partição de \(\Omega\) é uma coleção \(\{B_i : i \in I\}\) de eventos disjuntos
(no sentido de que \(B_i \in \mathcal{F}\) para cada \(i\), e \(B_i \cap B_j =
\varnothing\) se \(i \neq j\)) cuja união é \(\bigcup_i B_i = \Omega\). O
seguinte teorema da partição é extremamente útil.

\begin{theorem}[Teorema da Partição]\label{gw:teo:particao}
  Se \(\{B_1, B_2, \ldots\}\) é uma partição de \(\Omega\) com
  \(\mathbb{P}(B_i) > 0\) para cada \(i\), então
  \[
    \mathbb{P}(A) = \sum_i \mathbb{P}(A \mid B_i) \mathbb{P}(B_i) \quad
    \text{para } A \in \mathcal{F}.
  \]
\end{theorem}

Este teorema tem vários outros nomes pomposos, como ``teorema da probabilidade
total'', e está intimamente relacionado ao ``teorema de Bayes'',
\cref{gw:teo:bayes}.

\begin{proof}
  Temos que
  \[
    \begin{aligned}
      \mathbb{P}(A) &= \mathbb{P}\left(A \cap \left(\bigcup_i B_i\right)\right) \\
      &= \mathbb{P}\left(\bigcup_i (A \cap B_i)\right) \\
      &= \sum_i \mathbb{P}(A \cap B_i) \quad \text{pela aditividade
        enumerável de } \mathbb{P} \\
      &= \sum_i \mathbb{P}(A \mid B_i) \mathbb{P}(B_i) \quad \text{pela
        definição de probabilidade condicional, \eqref{eq:gw:prob-cond}}.
    \end{aligned}
  \]
\end{proof}

Aqui está um exemplo deste teorema em ação em um cálculo de dois estágios.

\begin{example}\label{gw:exemplo:atraso}
  Amanhã haverá chuva ou neve, mas não ambos; a probabilidade de chuva é
  \(\tfrac{2}{5}\) e a probabilidade de neve é \(\tfrac{3}{5}\). Se chover, a
  probabilidade de eu me atrasar para minha aula é \(\tfrac{1}{5}\), enquanto a
  probabilidade correspondente no caso de neve é \(\tfrac{3}{5}\). Qual é a
  probabilidade de eu me atrasar?
  \begin{solution}
    Seja \(A\) o evento de eu me atrasar e \(B\) o evento de chover. O par
    \(B, B^c\) é uma partição do espaço amostral (pois exatamente um deles
    deve ocorrer). Pelo \cref{gw:teo:particao},
    \[
      \mathbb{P}(A) = \mathbb{P}(A \mid B) \mathbb{P}(B) + \mathbb{P}(A \mid
      B^c) \mathbb{P}(B^c)
    \]
    \[
      = \frac{1}{5} \cdot \frac{2}{5} + \frac{3}{5} \cdot \frac{3}{5} =
      \frac{11}{25}.
    \]
  \end{solution}
\end{example}

Há muitas situações práticas em que se deseja deduzir algo a partir de uma
evidência. Escrevemos \(A\) para a evidência, e \(B_1, B_2, \ldots\) para os
possíveis ``estados da natureza''. Suponha que haja boas estimativas para as
probabilidades condicionais \(\mathbb{P}(A \mid B_i)\), mas que se busque, em
vez disso, uma probabilidade da forma \(\mathbb{P}(B_j \mid A)\).

\begin{theorem}[Teorema de Bayes]\label{gw:teo:bayes}
  Seja \(\{B_1, B_2, \ldots\}\) uma partição do espaço amostral \(\Omega\) tal
  que \(\mathbb{P}(B_i) > 0\) para cada \(i\). Para qualquer evento \(A\) com
  \(\mathbb{P}(A) > 0\),
  \[
    \mathbb{P}(B_j \mid A) = \frac{\mathbb{P}(A \mid
    B_j)\mathbb{P}(B_j)}{\sum_i \mathbb{P}(A \mid B_i)\mathbb{P}(B_i)}.
  \]
\end{theorem}

\begin{proof}
  Pela definição de probabilidade condicional (veja
  \cref{gw:exercicio:bayes-formula}),
  \[
    \mathbb{P}(B_j \mid A) = \frac{\mathbb{P}(A \mid
    B_j)\mathbb{P}(B_j)}{\mathbb{P}(A)},
  \]
  e a afirmação segue do teorema da partição, \cref{gw:teo:particao}.
\end{proof}

\begin{example}[Falsos positivos]\label{gw:exemplo:falsos-positivos}
  Uma doença rara, mas potencialmente fatal, tem incidência de 1 em \(10^5\) na
  população em geral. Há um teste diagnóstico, mas ele é imperfeito. Se você
  tem a doença, o teste dá positivo com probabilidade \(\tfrac{9}{10}\); se
  você não tem, o teste dá positivo com probabilidade \(\tfrac{1}{20}\). O
  resultado do seu teste é positivo. Qual é a probabilidade de você ter a
  doença?
  \begin{solution}
    Escreva \(D\) para o evento de você ter a doença, e \(P\) para o evento de
    o teste ser positivo. Pelo teorema de Bayes, \cref{gw:teo:bayes},
    \[
      \begin{aligned}
        \mathbb{P}(D \mid P) &= \frac{\mathbb{P}(P \mid
          D)\mathbb{P}(D)}{\mathbb{P}(P \mid D)\mathbb{P}(D) + \mathbb{P}(P
          \mid D^c)\mathbb{P}(D^c)} \\
        &= \frac{\frac{9}{10} \cdot \frac{1}{10^5}}{\frac{9}{10} \cdot
          \frac{1}{10^5} + \frac{1}{20} \cdot \frac{10^5 - 1}{10^5}}
          \approx 0.0002.
      \end{aligned}
    \]
  \end{solution}
  É mais provável que o resultado do teste esteja incorreto do que você tenha
  a doença.
\end{example}

\begin{exercise}\label{gw:exercicio:urnas}
  Aqui estão dois problemas rotineiros sobre bolas em urnas. Você recebe duas
  urnas. A urna I contém 3 bolas brancas e 4 pretas, e a urna II contém 2
  bolas brancas e 6 pretas.
  \begin{enumerate}[label=(\alph*)]
    \item Você retira uma bola ao acaso da urna I e a coloca na urna II. Em
          seguida, você retira uma bola ao acaso da urna II. Qual é a
          probabilidade de a bola ser preta?
    \item Desta vez, você escolhe uma urna ao acaso, cada uma das duas urnas
          sendo escolhida com probabilidade \(\tfrac{1}{3}\), e você retira uma
          % TODO-OCR[OCR-0002][JSON p.25]: o JSON lê "each of the two urns
          % being picked with probability 1/3"; mantido conforme a leitura
          % (decisão da Sessão 03). Verificar em uma passada visual se a
          % versão impressa mostra 1/3 ou 1/2.
          bola ao acaso da urna escolhida. Dado que a bola é preta, qual é a
          probabilidade de você ter escolhido a urna I?
  \end{enumerate}
\end{exercise}

\begin{exercise}\label{gw:exercicio:moeda-viciada}
  Uma moeda viciada mostra cara com probabilidade \(p = 1 - q\) a cada
  lançamento. Seja \(u_n\) a probabilidade de que, em \(n\) lançamentos, não
  ocorram duas caras consecutivas. Mostre que, para \(n \ge 1\),
  \[
    u_{n+2} = q u_{n+1} + p q u_n,
  \]
  e determine \(u_n\) pelo método usual (descrito no Apêndice B) quando
  \(p = \tfrac{2}{3}\).
\end{exercise}
